% =========================================================
% The Five Archetypes
% =========================================================
\paragraph{The Five Proof Archetypes.}

Every mathematical proof, at the highest level of compression, is an
instance of one of five archetypes. The statement map tells you the
specific opening move. The archetypes tell you the \emph{global shape}
of the argument --- what kind of reasoning is being done, and why it
constitutes a valid proof of the claim.

\begin{tcolorbox}[colback=propbox, colframe=propborder, arc=2pt,
  left=6pt, right=6pt, top=4pt, bottom=4pt,
  title={\small\textbf{The Five Proof Archetypes}}, fonttitle=\small\bfseries]
\begin{enumerate}
  \item \textbf{Construct something.}
    Exhibit a concrete object and verify it has every required property.
    This is the proof structure for all existence claims.

  \item \textbf{Assume two and compare.}
    Assume two objects both satisfy a definition; derive that they are
    equal. This is the proof structure for all uniqueness claims.

  \item \textbf{Take a minimal element.}
    Apply the well-ordering principle to a nonempty set of counterexamples;
    exploit the minimality of the least element to derive a contradiction.
    Used in number theory, descent arguments, and equivalence with induction.

  \item \textbf{Induct.}
    Verify a base case; show the property propagates to successors.
    Used whenever the claim concerns all natural numbers or any
    recursively defined structure.

  \item \textbf{Chase an arbitrary element.}
    Take an arbitrary element satisfying one condition; track it through
    definitions and established facts until it satisfies the target condition.
    Used for set inclusions, function properties, and direct proofs where
    the logic flows in a straight line from hypothesis to conclusion.
\end{enumerate}
\end{tcolorbox}

\begin{remark}[Why five archetypes, not more?]
The five archetypes are not a classification chosen for convenience.
They correspond to the five fundamental proof obligations that
mathematical statements impose:

To prove something \emph{exists}, you must produce it. To prove something
is \emph{unique}, you must show that any two instances are equal.
To prove a statement holds \emph{for all natural numbers}, the inductive
structure of $\mathbb{N}$ provides the only available tool. To prove a
\emph{conditional} claim $P \Rightarrow Q$, you must trace an arbitrary
element or object satisfying $P$ to a situation where $Q$ holds.
The minimal element archetype is the well-ordering counterpart to
induction, available whenever you need to reason about the smallest
counterexample.

Everything else in proof writing is refinement of these five. The specific
proof structures in Section~\ref{sec:proof-structures} (direct proof,
contrapositive, contradiction, cases, existence-uniqueness) are
instantiations of these archetypes adapted to specific logical forms.
The tactics in Section~\ref{sec:tactics} are the move-level tools
used to execute whichever archetype applies.
\end{remark}

\begin{remark}[Archetypes can nest, and this is a feature]
A single proof often uses more than one archetype. The most common
nested pattern in abstract algebra is: construct a witness (archetype~1),
then invoke a uniqueness theorem (archetype~2) to conclude that the
constructed object equals a previously named one.

This is the satisfy-and-cite tactic (Section~\ref{sec:satisfy-cite}).
It appears in almost every proof of the form ``$a = b$'' where $b$ is
an algebraically defined object. The nesting is not a complication ---
it is the standard pattern for proving equalities in algebraic structures,
and once recognized, it becomes mechanical.

Another common nesting: induct (archetype~4), and inside the inductive
step, use a direct proof argument (archetype~5) to move from $P(n)$ to
$P(n+1)$. The outer structure is induction; the inner argument is
element-chasing.
\end{remark}

\begin{remark}[The archetype identifies the proof's centre of gravity]
When you read a finished proof in a textbook, the archetype is usually
visible in the first line or two. ``Suppose $x, y$ both satisfy the
definition'' --- that is archetype~2. ``Define $f : X \to Y$ by
$f(x) = \ldots$'' --- that is archetype~1. ``Let $n \in \mathbb{N}$ be
arbitrary'' followed by induction markup --- that is archetype~4.
``Let $x \in A$'' followed by a chain of implications --- that is
archetype~5.

Reading proofs actively means identifying the archetype at the start
and then watching how the author executes it. This is far more efficient
than trying to absorb a proof line by line without a sense of its
overall shape.
\end{remark}
